%
\renewcommand{\theequation}{\theenumi}
\begin{enumerate}[label=\arabic*.,ref=\thesubsection.\theenumi]
\numberwithin{equation}{enumi}

\item	In Fig. \ref{fig:tri_med_centroid}, $BE$ and $CF$ are the medians. Show that $\triangle GEF \sim \triangle GBC$.
%
\begin{figure}[!ht]
	\begin{center}
		\resizebox{\columnwidth}{!}{\input{./figs/triangle/tri_med_centroid.tex}}
	\end{center}
	\caption{Similar Triangles}
	\label{fig:tri_med_centroid}	
\end{figure}
\\
\solution $\because EF \parallel BC$, alternate interior angles are equal.  Hence, corresponding angles of both triangles are equal and the triangles are similar.
\item Show that

%	Using Fig. \ref{ch2_median_ratio_val}, 
	\begin{align}
\label{eq:tri_med_centroid_ratio}
	\frac{GB}{GE} = \frac{GC}{GF} = 2
	\end{align}
%
\solution $\because \triangle GEF \sim \triangle GBC$ and $\frac{EF}{BC} =2$, \eqref{eq:tri_med_centroid_ratio} follows.
%
\item In Fig. \ref{fig:tri_med_meet}, AG is extended to meet $BC$ at $D$.	Show that 
	\begin{align}
\label{eq:tri_med_centroid_ratio_AD}
\frac{GA}{GD} = 2
	\end{align}
\begin{figure}[!ht]
	\begin{center}
		\resizebox{\columnwidth}{!}{\input{./figs/triangle/tri_med_meet.tex}}
	\end{center}
	\caption{Similar Triangles}
	\label{fig:tri_med_meet}	
\end{figure}

\solution  From Theorem \ref{them:tri_med_bpt}, $ EF || BC$.  Hence,
%
\begin{align}
\triangle APE &\sim \triangle ADC \quad (AAA)
\\
\implies AP &= PD = \frac{AD}{2}
\\
\implies  AG - GP &= GP + GD 
\\
\text{or, }  GP=\frac{AG - GD}{2} 
\label{eq:tri_med_centroid_ratio_AD1}
\end{align}
%
Similarly, 
%
\begin{align}
\triangle PGE &\sim \triangle BGD  \quad (AAA)
\\
\implies \frac{GP}{GD} &= \frac{GE}{GB} = \frac{1}{2}
\\
\text{or, }  GP=\frac{GD}{2} 
\label{eq:tri_med_centroid_ratio_AD2}
\end{align}
using \eqref{eq:tri_med_centroid_ratio}. From \label{eq:tri_med_centroid_ratio_AD1}
and \label{eq:tri_med_centroid_ratio_AD2},
%
\begin{align}
GP=\frac{AG - GD}{2} &=\frac{GD}{2}  
\\
\implies  \frac{AG}{GD} &=2
\end{align}
%%
\item Show that 
%
\begin{align}
BD = DC
\end{align}
%%
\item Show that the medians of a triangle meet at a point.
\\
\solution In Fig. \ref{fig:tri_med_meet}, $BE$ and $CF$	are medians that meet at $\vec{G}$.  Hence, 
from \eqref{tri_med_centroid_ratio}, 
%
\begin{align}
\label{eq:tri_centroid_g}
\vec{G} = \frac{2\vec{B}+\vec{E}}{3}
\end{align}
%
If the median $AD$ meets $BE$ in $\vec{G}_1$, from Theorem \eqref{tri_med_centroid_ratio}, 
%
\begin{align}
\vec{G}_1 = \frac{2\vec{B}+\vec{E}}{3}
\label{eq:tri_centroid_g1}
\end{align}
%
From  \eqref{eq:tri_centroid_g}
 and \eqref{eq:tri_centroid_g1}, $\vec{G} = \vec{G}_1$.  Thus, the medians meet at $\vec{G}$.
\item The centroid divides the median in the ratio $2:1$.
\label{them:tri_med_ratio}



%\item
%	In Fig. \ref{ch2_median_3}, the line $AO$ cuts $EF$ at $G$ and is extended to meet the side $BC$ at $D$.  Show that 
%	%
%	\begin{equation}
%	\frac{OA}{OD} = 2.
%	\end{equation}
%	%  
%
%\begin{figure}[!ht]
%	\begin{center}
%		
%		%\includegraphics[width=\columnwidth]{./figs/ch2_median_3}
%		%\vspace*{-10cm}
%		\resizebox{\columnwidth}{!}{\input{./figs/fig_2.14.tex}}
%	\end{center}
%	\caption{Similar Triangles and Median}
%	\label{ch2_median_3}	
%\end{figure}
%
%
%\item
%	In Fig. \ref{ch2_median_ratio_val}, show that
%	\begin{equation}
%	ar\brak{\triangle AFE} = 	\frac{1}{4} ar\brak{\triangle ABC}
%	\end{equation}
%	\label{ch2_small_triang}
%%
%\solution We have
%%%
%\begin{align}
%ar\brak{\triangle AFE} &= \frac{1}{2} \frac{b}{2}\frac{c}{2}\sin A = \frac{1}{4}.\frac{1}{2}bc \sin A \\
%&=\frac{1}{4}ar\brak{\triangle ABC}
%\end{align}
%%
%		%\includegraphics[width=\columnwidth]{./figs/ch2_median_ratio_val}
%		%\vspace*{-10cm}
%
%%
%\begin{figure}[!ht]
%	\begin{center}
%		
%		\resizebox{\columnwidth}{!}{\input{./figs/fig_2.9.tex}}
%	\end{center}
%	\caption{$O$ divides medians in the ratio $2:1$}
%	\label{ch2_median_ratio_val}	
%\end{figure}
%%
%
%\solution Using the sine formula and \eqref{ch2_small_triang}, areas of some triangles in Fig. \ref{ch2_median_ratio_val} are listed in the following table
%\begin{center}
%\begin{tabular}{|c|c|}
%	\hline  Triangle& Area  \\ 
%	\hline  OFE & $\frac{1}{2}pq \sin \theta $ \\ 
%	\hline  BOF &  $\frac{k}{2}pq \sin \theta$\\ 
%	\hline  COE & $\frac{k}{2}pq \sin \theta$ \\ 
%	\hline  BOC&  $\frac{k^2}{2}pq \sin \theta$\\
%	\hline  BOC&  $\frac{1}{2} ar \brak{\triangle ABC}$\\  
%	\hline 
%\end{tabular} 
%\end{center}
%%
%where we have used the fact that 
%
%\begin{equation}
%\sin \angle BOC = \sin \brak{180^{\degree}-\theta} = \sin \theta
%\end{equation}
%%
%Since $BE$ is the median
%\begin{align}
%\begin{split}
%ar\brak{\triangle BEC} &= \frac{1}{2} ar\brak{\triangle ABC}  
%\\
%& = ar\brak{\triangle BOC} + ar\brak{\triangle COE}\\
%ar\brak{\triangle BEA} &= \frac{1}{2} ar\brak{\triangle ABC}
%\\
%& = ar\brak{\triangle AFE} + ar\brak{\triangle BOF} +  ar\brak{\triangle FOE}  
%\end{split}
%\end{align}
%%
%Substituting the respective values from the above table,
%\begin{align}
%\begin{split}
%\frac{1}{2} ar\brak{\triangle ABC} &= \frac{k^2}{2}pq \sin \theta + \frac{k}{2}pq \sin \theta
%\\
%\frac{1}{2} ar\brak{\triangle ABC} &= \frac{k}{2}pq \sin \theta + \frac{1}{2}pq \sin \theta + \frac{1}{4} ar\brak{\triangle ABC}
%\end{split}
%\end{align}
%%
%Simplifying the above,
%%
%\begin{align}
%k(k+1) = 2(k+1)
%\end{align}
%%
%Since $k \neq -1$, $k = 2$ and the proof is complete.
%

\end{enumerate}
